\section{A stable bundle with cubic deformation equations}\label{sec:bicubic}

We construct a compact locally free example for which the full obstruction
in Theorem~\ref{thm:main} is essential. The method is the
curve-to-bundle construction of Thomas \cite{thomas}, with a different
bicubic and a triple-line fiber.

\subsection{The compact geometry and the Hilbert germ}
On $B=\PP^2$ choose smooth cubic curves defined by $R,P,Q$, with
pairwise transverse intersections and no common point. Such triples
form a nonempty open set: choose transverse smooth $P,Q$, then a
general smooth $R$ avoiding their finite intersection and transverse
to both. Choose $p\in\{P=Q=0\}$, so $R(p)\ne0$. In
$\mathcal P=\PP^2_{[x:y:z]}\times B$ put
\begin{equation}\label{eq:bicubic}
 X=\{x^3R+y^3P+z^3Q=0\},\qquad C=\{x=0\}\times\{p\}.
\end{equation}
Write $\omega_1,\omega_2$ for the restrictions of the two hyperplane
classes.

\begin{lemma}\label{lem:bicubic-smooth}
The threefold $X$ is smooth and connected, $K_X\simeq\OO_X$,
$H^1(X,\OO_X)=H^2(X,\OO_X)=0$, and
$\operatorname{Pic}X=\Z\omega_1\oplus\Z\omega_2$.
The fiber over $p$ is the triple line $x^3=0$.
\end{lemma}

\begin{proof}
At a singular point the fiber derivatives imply
$x^2R=y^2P=z^2Q=0$. If all three fiber coordinates are nonzero,
all three coefficient cubics vanish, a contradiction.
If only one is nonzero, the corresponding smooth cubic has a
nonzero base differential. If two are nonzero, the corresponding
two cubics vanish transversely, and the nonzero coefficients
multiplying their independent differentials cannot cancel.
In both cases the total defining differential is nonzero.
Thus $X$ is smooth. Adjunction gives $K_X\simeq\OO_X$.
The hypersurface is ample, so Lefschetz gives the Picard assertion.
The restriction sequence of the bicubic gives connectedness and
the stated structure-sheaf cohomology. The fiber assertion follows
from $P(p)=Q(p)=0$ and $R(p)\ne0$.
\end{proof}

Near $p$, the functions $U=Q/R$ and $V=P/R$ are analytic coordinates.
The local equation is exactly
\begin{equation}\label{eq:local-bicubic}
 x^3+Vy^3+Uz^3=0.
\end{equation}

\begin{proposition}\label{prop:hilbert}
The normal bundle of $C$ in $X$ is
$\OO_C(1)\oplus\OO_C(-3)$, and its analytic Hilbert germ is
\begin{equation}\label{eq:hilbert-ring}
 \Specan\C\{a,b\}/(a^2b,ab^2).
\end{equation}
Its tangent dimension is two. Every first-order direction extends
to second order, whereas a direction with both coordinates nonzero
is obstructed at order three.
\end{proposition}

\begin{proof}
The ambient normal bundle is
$\OO_C(1)\oplus\OO_C\oplus\OO_C$.
The normal map to $\OO_C(3)$, obtained from
\eqref{eq:local-bicubic}, is
\[
 (s,c,d)\longmapsto cz^3+dy^3.
\]
It is surjective, with kernel $\OO_C(1)\oplus\OO_C(-3)$.
Since $H^1(N_{C/\mathcal P})=0$, the ambient Hilbert scheme is
smooth of dimension four at $C$. The natural family of lines in
nearby fibers has tangent isomorphism onto $H^0(N_{C/\mathcal P})$,
so it gives the full analytic Hilbert chart, including infinitesimal
deformations:
\[
 x=ay+bz,\qquad U=c,\qquad V=d.
\]
Substitution in \eqref{eq:local-bicubic} gives the four coefficients
\begin{equation}\label{eq:four-coefficients}
 a^3+d,\qquad 3a^2b,\qquad 3ab^2,\qquad b^3+c
\end{equation}
of $y^3,y^2z,yz^2,z^3$. Their vanishing is the exact relative
Hilbert condition. Eliminate $d=-a^3$ and $c=-b^3$ to obtain
\eqref{eq:hilbert-ring}.

The ideal has primary decomposition
\begin{equation}\label{eq:primary}
 (a^2b,ab^2)=(a)\cap(b)\cap(a,b)^3.
\end{equation}
Its radical is $(ab)$, and its nonzero minimal defining ideal lies
in the cube of the maximal ideal. The embedding dimension is two.
The order of this ideal is invariant under analytic changes of
minimal coordinates, so the full germ cannot be defined by
quadratic equations. In particular $a=b=\tau$ defines a deformation
over $\C[\tau]/(\tau^3)$ but cannot extend to
$\C[\tau]/(\tau^4)$ with the same first derivative: higher
coefficients of $a,b$ cannot cancel the leading cubic products.
\end{proof}

The reduced germ consists of two transverse smooth axes, with an
embedded component at their intersection in the full analytic germ.
Higher-order curve obstructions occur in other compact Calabi--Yau
examples; for instance \cite[Section~5.2]{ashok} computes fifth-order
initial equations for lines on the Fermat quintic. Here the cubic
germ will be realized by a stable locally free bundle.

\subsection{The Serre bundle and its stability}
Set $L=\OO_X(-2,2)$. Its restriction to $C$ equals
$\det N_{C/X}=\OO_C(-2)$. The ambient restriction sequence
\begin{equation}\label{eq:restriction}
 0\longrightarrow\OO_{\mathcal P}(i-3,j-3)
 \longrightarrow\OO_{\mathcal P}(i,j)
 \longrightarrow\OO_X(i,j)\longrightarrow0
\end{equation}
shows that both $L$ and $L^*$ are acyclic.
For $L$, the ambient bundles are $\OO_{\mathcal P}(-2,2)$ and
$\OO_{\mathcal P}(-5,-1)$; for $L^*$ they are
$\OO_{\mathcal P}(2,-2)$ and $\OO_{\mathcal P}(-1,-5)$.
Each is acyclic by the line-bundle cohomology of $\PP^2$.

The Hartshorne--Serre construction \cite{hartshorne,thomas} gives
a rank-two locally free bundle with a section:
\begin{equation}\label{eq:serre}
 0\longrightarrow\OO_X\xrightarrow{s}E
       \longrightarrow I_CL\longrightarrow0 .
\end{equation}
Indeed the local extension generator is a nowhere-zero section of
$\det N_{C/X}\otimes L^*|_C=\OO_C$, and
$H^1(L^*)=H^2(L^*)=0$ gives its unique global lift, up to the
usual scalar choice. This generator ensures local freeness along
$C$; away from $C$ one has an extension of line bundles.
Put
\[
 A=E(1,-1),\qquad \det A=\OO_X .
\]
Since $L|_C=\OO_C(-2)$, \eqref{eq:serre} and the restriction
sequence give
\begin{equation}\label{eq:section-cohomology}
 h^0(E)=1,\qquad H^1(E)=0 .
\end{equation}

\begin{proposition}\label{prop:stability}
The bundles $E$ and $A$ are slope-stable for
\[
 H_N=N\omega_1+\omega_2,\qquad 1<N<1+\sqrt3.
\]
In particular $A$ is stable of slope zero for the integral ample
class $H=2\omega_1+\omega_2$.
\end{proposition}

\begin{proof}
The double dual of a saturated rank-one subsheaf of $E$ is a line
bundle with the same first Chern class, and its map to the locally
free bundle extends across codimension two. By the Picard
description, it is $\OO_X(-k,-l)$ for some integers $k,l$.
It therefore suffices to show that every twist $E(k,l)$ with a
nonzero section has positive slope.

For all integers $i,j$ the ambient group
$H^1(\OO_{\mathcal P}(i-3,j-3))$ vanishes. Thus
\[
 H^0(\OO_X(i,j))\ne0\quad\Longleftrightarrow\quad i,j\ge0 .
\]
Moreover
\[
 H^0(I_C(i,j))\ne0
 \quad\Longleftrightarrow\quad i,j\ge0,\quad i+j\ge1.
\]
For $i\ge1$, use a polynomial divisible by $x$; for $j\ge1$,
use a base polynomial vanishing at $p$. Constants do not vanish on
$C$. Twisting \eqref{eq:serre} shows that a section of $E(k,l)$
can occur only in the union of the coordinatewise upper quadrants
based at
\[
 (0,0),\qquad (2,-1),\qquad (3,-2).
\]
The mixed triple intersections are
$\omega_1^2\omega_2=\omega_1\omega_2^2=3$, and the pure cubes
vanish. Writing degrees here without the common factor $1/2!$,
one obtains
\begin{equation}\label{eq:three-degrees}
 \begin{aligned}
 c_1(E)\cdot H_N^2&=6(N^2-1),\\
 c_1(E(2,-1))\cdot H_N^2&=6(2N+1),\\
 c_1(E(3,-2))\cdot H_N^2&=6(2+2N-N^2).
 \end{aligned}
\end{equation}
All are positive on the stated interval.
Adding a nonnegative twist increases degree since $H_N$ is ample.
Thus $\mu_{H_N}(E(k,l))>0$ whenever it has a section, equivalently
$\mu_{H_N}(\OO_X(-k,-l))<\mu_{H_N}(E)$.
This proves stability. Twisting preserves stability and gives
the result for $A$.
\end{proof}

The third quadrant is necessary for this proof. It accounts for
sections of $I_C(i,0)$ with $i\ge1$.
The argument proves the sufficient interval above independently of
any broader stability assertion for a different bicubic.

\subsection{Equality of the bundle and curve deformation germs}

\begin{theorem}\label{thm:cubic-bundle}
The fixed-$X$, fixed-determinant analytic deformation germ of $A$
is \eqref{eq:hilbert-ring}, including its nonreduced structure.
In particular $\dim\Ext^1(A,A)=2$, the quadratic Kuranishi
obstruction vanishes, and the cubic obstruction is nonzero in
minimal coordinates. Each reduced axis gives a genuine
one-parameter family of pairwise nonisomorphic stable locally
free bundles near $A$.
\end{theorem}

\begin{proof}
We compare the deformation functors, not only their tangent
dimensions. Let $T$ be a local Artin $\C$-algebra.
Every deformation of $E$ over $X\times\operatorname{Spec}T$
has a section lifting $s$: induction over small extensions
places the section-lifting obstruction in $H^1(E)$ tensored
with the square-zero kernel, which vanishes by
\eqref{eq:section-cohomology}. The same induction and $h^0(E)=1$
show that its module of sections is free of rank one.
The lifted section is unique up to a unit of $T$.
Since its closed-fiber zero is a regular codimension-two zero,
its zero scheme is a flat local complete intersection deformation
of $C$. Its determinant remains $L$, as
$H^1(X,\OO_X)=0$ and the determinant is fixed.

Conversely, every embedded deformation $C_T$ belongs to the ambient
line chart used in Proposition~\ref{prop:hilbert}. It is relatively
$\PP^1$, and its relative normal determinant equals $L|_{C_T}$
by adjunction: $K_X$ is trivial and
$\OO_{\mathcal P}(-2,2)|_{C_T}$ has relative degree $-2$.
The possible line from the local base is trivial.
The fixed acyclicity of $L^*$ persists after tensoring with $T$.
The relative Serre construction therefore extends the local normal
generator uniquely to a vector bundle $E_T$ with regular section.
Changing that generator by a unit only rescales the section and
gives an isomorphic bundle. These two constructions are inverse
on isomorphism classes and commute with base change.

Twisting by $\OO_X(1,-1)$ preserves the deformation functor.
Stability makes $A$ simple; scalar automorphisms act trivially on
its deformation base. The completed local ring is therefore the
completion of \eqref{eq:hilbert-ring}.
The constructions also apply to small analytic families:
vanishing of $H^1(E)$, the one-dimensional space of sections,
analytic base change and regularity of the zero scheme give
inverse analytic germ maps. Thus the assertion includes the
convergent structure and nilpotents, not only the completed ring.

The defining ideal has order three in minimal embedding dimension
two. Hence a minimal Kuranishi map has zero quadratic part and
nonzero cubic part. The two reduced axes give distinct points of
the simple-bundle deformation germ. Stability and local freeness
persist in a small neighborhood.
\end{proof}

Thomas's corrected example \cite[Addendum]{thomas} has ring
$\C[\varepsilon,\eta]/(\varepsilon^2,\eta^2)$.
The triple-line geometry \eqref{eq:bicubic} instead gives a
one-dimensional germ with cubic initial ideal. The use of the
Serre method is common to both constructions.

\subsection{The full repeated-factor application}
Choose a K\"ahler form in $H=2\omega_1+\omega_2$ and the normalized
HE metric on $A$. Set $E_0=A\otimes\C^3$.
Then
\[
 V=\Ext^1(A,A)\otimes\End\C^3,\qquad \dim_\C V=18,\qquad N=8.
\]
Since $H^1(X,\OO_X)=0$, the determinant condition removes no
degree-one harmonic direction.

\begin{corollary}\label{cor:full-application}
The full local HYM problem at $A\otimes\C^3$, with all eighteen
harmonic coordinates and the complete convergent matrix obstruction,
satisfies every conclusion of Theorem~\ref{thm:main}.
At fixed polarization the relative error is $O(\|x\|^2)$.
There are eight small trace-free endomorphism eigenvalues, including
their exact changing kernel, a uniform complementary gap, and the
same relative harmonic metric comparison on genuine smooth pieces.
Every sufficiently nearby actual solution is represented.
\end{corollary}

\begin{proof}
Lemmas and propositions above verify compact smooth K\"ahler
geometry, local freeness, stability, determinant and the analytic
deformation germ. Thus all background hypotheses of
Theorem~\ref{thm:main} hold. Ext additivity gives the eighteen
coordinates and the full multiplicity algebra gives the eight
trace-free zero-form directions. Apply the theorem to this entire
family. Small polarization variations are covered as well, by the
theorem or by Proposition~\ref{prop:moving} in its moving coordinates.
\end{proof}

The scalar slice $u\otimes I_3$ retains the cubic obstruction of
$A$ by the blockwise canonical recursion. Its quadratic obstruction
is zero, so a quadratic truncation would admit false actual
deformations. No claim identifies the complete matrix relations
by substituting noncommuting matrices into the two scalar equations.

There is an explicit spectral family with arbitrary scale
separation. Parametrize one reduced axis in canonical harmonic
coordinates by
\[
 f(z)=ze+h(z),\qquad \|e\|=1,\quad h(z)\perp e,\quad h(z)=O(z^2).
\]
Use the harmonic linear coordinate along $e$ as parameter.
Holomorphic divided differences give
\begin{equation}\label{eq:divided-difference}
 \|f(z)-f(w)\|^2
   =|z-w|^2\bigl(1+O((|z|+|w|)^2)\bigr).
\end{equation}
Indeed $\|h(z)-h(w)\|\le C(|z|+|w|)|z-w|$, and the two
summands are orthogonal. The estimate remains relative when
$z-w$ is arbitrarily small.

For any three nearby parameters put
$x=\diag(f(z_1),f(z_2),f(z_3))$.
The harmonic reconstruction is blockwise, hence integrable.
Every loop coefficient is diagonal and normal, so $\mu(x)=0$.
The bundle is $A_{z_1}\oplus A_{z_2}\oplus A_{z_3}$.
On each off-diagonal matrix unit the mass is $2\Vol(\omega_0)$
and the commutator energy is $\|f(z_i)-f(z_j)\|^2$.
Thus for every unordered pair the finite eigenvalue is
\begin{equation}\label{eq:pair-spectrum}
 \nu_{ij}=\frac{\|f(z_i)-f(z_j)\|^2}{2\Vol(\omega_0)}
\end{equation}
with complex multiplicity two. Distinct parameters give trace-free
kernel dimension two; one double collision gives dimension four;
a triple collision gives dimension eight.
The corresponding compact stabilizers before removing the
ineffective scalar are $\mathrm U(1)^3$,
$\mathrm U(1)\times\mathrm U(2)$ and $\mathrm U(3)$.

Take $(z_1,z_2,z_3)=(0,s,s+\varepsilon)$, where
$0<\varepsilon\le s\to0$. Equations
\eqref{eq:divided-difference}--\eqref{eq:pair-spectrum} and
Theorem~\ref{thm:main} give two exact trace-free zeros and six
positive eigenvalues, with two entries at each scale
\begin{equation}\label{eq:separated-spectrum}
 \frac{\varepsilon^2}{2\Vol(\omega_0)}(1+O(s^2)),\quad
 \frac{s^2}{2\Vol(\omega_0)}(1+O(s^2)),\quad
 \frac{(s+\varepsilon)^2}{2\Vol(\omega_0)}(1+O(s^2)).
\end{equation}
This means comparison of ordered lists with repetitions, not
exact equality of perturbed pairs. The constant is independent
of $\varepsilon/s$, and the ninth trace-free eigenvalue has a
fixed positive lower bound. The kernel dimensions at collisions
are exact. The quotient metric conclusion holds on the full local
family, with additional analytic refinement at the crossing of
the reduced axes. The smooth-branch calculation in
\eqref{eq:separated-spectrum} tests relative uniformity; it does
not assign a spectral exponent to the nilpotent structure.

